\begin{resumo}[Abstract]
  \begin{otherlanguage*}{english}
    Software quality is a central concern in Software Engineering, with maintainability standing out as one of its most relevant characteristics for the long-term evolution and sustainability of systems. Traditional static analysis tools, however, present limitations in capturing semantic and contextual aspects of source code, a gap that motivates the investigation of Large Language Models (LLMs) as a complementary evaluation approach. Within this scenario, the way instructions are formulated — prompt engineering patterns — emerges as a potentially relevant factor for the results produced by these models. This work aims to investigate how different prompt engineering patterns influence the assessment of the maintainability indicator of software products by LLMs, taking established software quality models as reference. The research was structured using the Goal-Question-Metric (GQM) approach to define the research question, followed by a structured literature review, conducted with support from the PICO protocol on the Scopus database, resulting in 38 analyzed studies. The synthesis of findings showed that maintainability is predominantly operationalized through measures such as complexity, code smells, duplication, and technical debt. Additionally, it was found that there is no specific prompt pattern consolidated for maintainability, although strategies that reduce the analysis context window and divide requests into smaller tasks tend to favor LLM performance. Based on this evidence, a controlled experiment was planned, structured according to a Randomized Complete Block Design (RCBD), in which three prompt patterns (zero-shot, few-shot, and chain-of-thought) will be applied to source code files from open-source projects in the web application domain, using three LLM tools (ChatGPT, Claude, and DeepSeek), with two reference quality models (Q-Rapids and MeasureSoftGram) serving as block variables, and the SQALE method (via SonarQube) being used for complementary triangulation of the results. By executing this experiment in the subsequent stage of this work, it is expected to verify whether significant differences exist between prompt patterns in the assessment of maintainability, and to what extent the results obtained by LLMs align with quality models traditionally used in Software Engineering.

    \vspace{\onelineskip}
  
    \noindent 
    \textbf{Key-words}: Large Language Models. Software Maintainability. Prompt Engineering. Controlled Experiment. Software Quality.
  \end{otherlanguage*}
\end{resumo}
